\documentclass[12pt,a4paper]{article}
\usepackage[utf8]{inputenc}
\usepackage[french]{babel}
\usepackage{amsmath,amsfonts,amssymb,mathrsfs}
\usepackage{geometry}
\geometry{hmargin=2cm,vmargin=2cm}

\begin{document}

\begin{center}
    \large\textbf{Composition du 2\textsuperscript{e} Semestre -- Classe : 1S2}
\end{center}

\hrulefill

\subsection*{Exercice 1 \hfill (5 points)}

\noindent\textbf{Recopier puis compléter les phrases suivantes : (2 pts)}
\begin{enumerate}
    \item[a)] La droite $(D)$ d'équation $x=2$ est un axe de symétrie à la courbe d'une fonction $f$ lorsque \dots
    \item[b)] Une fonction $f$ est continue sur un intervalle $I$ lorsque \dots
    \item[c)] Une fonction $f$ est prolongeable par continuité au point d'abscisse $x_{0}$ lorsque \dots
\end{enumerate}

\noindent\textbf{Pour chacune des affirmations suivantes, dire si c'est vrai ou faux en justifiant.}
\begin{enumerate}
    \item[\textbf{1.}] Soit $f$ une fonction définie par : $f(x)=\dfrac{x^{3}-3x^{2}+x-5}{x-x^{2}}$ \\
    $\lim_{x\rightarrow-\infty}\left(1+\dfrac{f(x)}{x}\right)=2$ \hfill (1 pt)
    
    \item[\textbf{2.}] La fonction $f$ définie par $f(x)=|x|$ est dérivable en 0. \hfill (1 pt)
    
    \item[\textbf{3.}] Soit $f$ une fonction définie par $f(x)=\dfrac{(1-3x)^{2}}{2x-1}$, alors $f^{\prime}(x)=\dfrac{2(1-3x)(2-3x)}{(2x-1)^{2}}$. \hfill (1 pt)
\end{enumerate}

\smallbreak
\hrulefill
\smallbreak

\subsection*{Exercice 2 \hfill (5 points)}

Soit $ABCD$ un quadrilatère quelconque, $I$ et $J$ les milieux respectifs de $[AB]$ et $[BC]$.\\
On note $G$ le centre de gravité du triangle $ABC$.\\
On considère les points $L$ et $K$ tels que :
\[ L=\text{bar}\{(A,1);(D,3)\} \quad \text{et} \quad K=\text{bar}\{(C,1);(D,3)\} \]

\begin{enumerate}
    \item Faire une figure et placer les points $G$, $K$ et $L$.
    \item Soit $H=\text{bar}\{(A,1); (B, 1); (C, 1); (D, 3)\}$
    \begin{enumerate}
        \item Montrer que $H$ est le barycentre des points $G$ et $D$ affectés des coefficients à préciser.
        \item En déduire que $H \in (GD)$.
        \item Démontrer que les droites $(GD)$, $(JL)$ et $(IK)$ sont concourantes en un point que l'on précisera.
    \end{enumerate}
    \item Déterminer puis construire l'ensemble des points $M$ du plan tels que :
    \begin{enumerate}
        \item $(\mathcal{E}) : \|\overrightarrow{MA}+\overrightarrow{MB}+\overrightarrow{MC}\|=9$
        \item $(\mathcal{F}) : \|\overrightarrow{MA}+\overrightarrow{MB}+\overrightarrow{MC}+3\overrightarrow{MD}\|=\dfrac{3}{2}\|\overrightarrow{MA}+3\overrightarrow{MD}\|$
    \end{enumerate}
\end{enumerate}

\smallbreak
\hrulefill
\smallbreak

\subsection*{Exercice 3 \hfill (10 points)}

\begin{enumerate}
    \item \textbf{Calculer les limites suivantes :}
    \begin{enumerate}
        \item $\lim_{x\rightarrow-\frac{2}{3}^{+}}\dfrac{-x^{2}+5x+1}{3x+2}$
        \item $\lim_{x\rightarrow-\infty}\sqrt{x^{2}+3x}+3x$
        \item $\lim_{x\rightarrow0}\dfrac{\sqrt{x+1}-1}{\sqrt{x+9}-3}$
    \end{enumerate}
    
    \item \textbf{Soit $f$ la fonction définie par :}
    \[ f(x)=\begin{cases}
    \sqrt{|x^{2}-x|} & \text{si } x\ge0 \\ 
    \dfrac{x^{2}}{x^{2}-1} & \text{si } x<0
    \end{cases} \]
    
    \begin{enumerate}
        \item Déterminer le domaine de définition de $f$.
        \item Calculer les limites aux bornes de $D_{f}$ puis préciser les asymptotes éventuelles.
        \item Étudier la continuité de $f$ en 0.
        \item Étudier la continuité de $f$ sur son ensemble de définition.
        \item Écrire $f$ sans symbole de la valeur absolue.
        \item Montrer que la droite $(D)$ d'équation $y=x-\dfrac{1}{2}$ est asymptote oblique à $(C_{f})$ en $+\infty$.
        \item Étudier le sens de variation puis dresser le tableau de variation de $f(x)$.
    \end{enumerate}
\end{enumerate}

\begin{center}
    \textbf{Bonne chance!}
\end{center}

\end{document}